\subsection{Resultados para Grover}

Grover fue el circuito más favorable para la compresión del estado. En los seis escenarios
evaluados, Blosc produjo los mayores ahorros de memoria y lo hizo con un patrón muy estable a
medida que aumentó el número de qubits. ZFP introdujo errores máximos absolutos del orden de
$10^{-14}$, pero aun admitiendo esa pérdida controlada redujo menos memoria y presentó una
penalización temporal mucho mayor. En consecuencia, la principal conclusión
de esta subsección es que, dentro de la campaña experimental realizada, la alternativa más
coherente para Grover es Blosc cuando la restricción dominante es la memoria disponible.

\subsubsection{Objetivo único}

La Tabla~\ref{tab:grover-single-results} resume los promedios de las tres posiciones
evaluadas para el estado marcado: $N/8$, $N/2$ y $7N/8$. Las variaciones entre estas tres
ubicaciones fueron pequeñas en cada tamaño, por lo que el promedio preserva adecuadamente la
tendencia general del caso de objetivo único.

\begin{table}[H]
    \centering
    \small
    \setlength{\tabcolsep}{3pt}
    \caption{Grover con objetivo único: promedios sobre las posiciones $N/8$, $N/2$ y $7N/8$}
    \label{tab:grover-single-results}
    \resizebox{\linewidth}{!}{%
    \begin{tabular}{lcccccc}
        \hline
        \textbf{Qubits} & \textbf{Estrategia} & \textbf{Mem. pico (MB)} & \textbf{Ahorro mem. (\%)} & \textbf{Tiempo (s)} & \textbf{Tiempo rel. (x)} & \textbf{Error} \\
        \hline
        22 & Referencia & 71.6 & 0.0 & 7.66 & 1.00 & 0 \\
        22 & Blosc & 15.9 & 77.8 & 37.16 & 4.85 & 0 \\
        22 & ZFP & 45.5 & 36.4 & 289.30 & 37.78 & $1.61\times 10^{-14}$ \\
        23 & Referencia & 135.5 & 0.0 & 24.82 & 1.00 & 0 \\
        23 & Blosc & 16.2 & 88.1 & 104.85 & 4.22 & 0 \\
        23 & ZFP & 78.2 & 42.3 & 801.01 & 32.27 & $1.11\times 10^{-14}$ \\
        24 & Referencia & 263.7 & 0.0 & 65.40 & 1.00 & 0 \\
        24 & Blosc & 17.1 & 93.5 & 279.33 & 4.27 & 0 \\
        24 & ZFP & 143.1 & 45.7 & 2274.03 & 34.77 & $8.04\times 10^{-15}$ \\
        \hline
    \end{tabular}
    }
\end{table}

El patrón del caso de objetivo único es claro. Blosc ahorra entre 77.8\% y 93.5\% de memoria
frente a la referencia, y el ahorro aumenta con el número de qubits. La penalización temporal
permanece relativamente estable, entre $4.22\times$ y $4.85\times$. ZFP también reduce la
memoria, pero solo entre 36.4\% y 45.7\%, y lo hace con un costo temporal muy superior, entre
$32.27\times$ y $37.78\times$. Además, el error máximo absoluto de ZFP se mantiene entre
$8.04\times10^{-15}$ y $1.61\times10^{-14}$, es decir, en un rango muy pequeño. Sin embargo,
incluso admitiendo ese error, ZFP no supera a Blosc ni en memoria ni en tiempo. Esto indica
que, para Grover con objetivo único, la compresión sin pérdida ofrece la relación más
favorable entre capacidad de ahorro y costo adicional.

\subsubsection{Multiobjetivo}

La Tabla~\ref{tab:grover-multi-results} reporta los resultados del caso multiobjetivo. Aunque
la presencia de varios estados marcados reduce el número de iteraciones de Grover y con ello
el tiempo absoluto del circuito, el patrón relativo entre estrategias permanece muy próximo al
observado en el caso de objetivo único.

\begin{table}[H]
    \centering
    \small
    \setlength{\tabcolsep}{3pt}
    \caption{Grover multiobjetivo con tres estados marcados}
    \label{tab:grover-multi-results}
    \resizebox{\linewidth}{!}{%
    \begin{tabular}{lcccccc}
        \hline
        \textbf{Qubits} & \textbf{Estrategia} & \textbf{Mem. pico (MB)} & \textbf{Ahorro mem. (\%)} & \textbf{Tiempo (s)} & \textbf{Tiempo rel. (x)} & \textbf{Error} \\
        \hline
        22 & Referencia & 71.6 & 0.0 & 4.46 & 1.00 & 0 \\
        22 & Blosc & 15.9 & 77.8 & 21.68 & 4.86 & 0 \\
        22 & ZFP & 45.3 & 36.7 & 170.16 & 38.13 & $1.61\times 10^{-14}$ \\
        23 & Referencia & 135.6 & 0.0 & 13.77 & 1.00 & 0 \\
        23 & Blosc & 16.2 & 88.0 & 60.54 & 4.40 & 0 \\
        23 & ZFP & 78.3 & 42.2 & 468.60 & 34.02 & $1.11\times 10^{-14}$ \\
        24 & Referencia & 263.7 & 0.0 & 40.53 & 1.00 & 0 \\
        24 & Blosc & 17.3 & 93.4 & 168.84 & 4.17 & 0 \\
        24 & ZFP & 143.3 & 45.7 & 1342.57 & 33.12 & $8.04\times 10^{-15}$ \\
        \hline
    \end{tabular}
    }
\end{table}

La lectura del caso multiobjetivo confirma la conclusión anterior. Blosc conserva un ahorro de
memoria entre 77.8\% y 93.4\%, prácticamente idéntico al del objetivo único, y mantiene una
penalización temporal cercana a $4.2\times$--$4.9\times$. ZFP vuelve a situarse en un punto
intermedio de ahorro espacial, entre 36.7\% y 45.7\%, con sobrecostos temporales que superan
los $33\times$ en todos los tamaños y errores máximos absolutos nuevamente del orden de
$10^{-14}$. La diferencia principal frente al caso de objetivo único no es cualitativa, sino
cuantitativa: la referencia no comprimida requiere menos iteraciones y por ello el tiempo
absoluto del circuito disminuye. Desde el punto de vista comparativo, el resultado sigue
siendo el mismo: la pequeña discrepancia numérica admitida por ZFP no se traduce en una
ventaja práctica frente a Blosc.

\subsubsection{Comparación complementaria con la variante optimizada}

La variante de parametrización reducida descrita en la
Sección~\ref{subsec:grover_reducido_impl} se evaluó como referencia complementaria para
estimar el efecto de explotar explícitamente la simetría global de Grover. En esta
subsección se mantiene el mismo criterio de lectura usado en el resto del capítulo: el ahorro
de memoria y el tiempo relativo se calculan frente a la referencia no comprimida de la propia
variante optimizada, para el mismo caso y el mismo número de qubits. En el caso de objetivo
único, cada fila corresponde al promedio de las tres posiciones evaluadas para ese tamaño:
$N/8$, $N/2$ y $7N/8$. Dado que el interés aquí es exclusivamente espacial y temporal, la
tabla omite la columna de error.

\begin{table}[H]
    \centering
    \small
    \setlength{\tabcolsep}{3pt}
    \caption{Variante optimizada de Grover con objetivo único}
    \label{tab:grover-optimized-single}
    \resizebox{\linewidth}{!}{%
    \begin{tabular}{lccccc}
        \hline
        \textbf{Qubits} & \textbf{Estrategia} & \textbf{Mem. pico (MB)} & \textbf{Ahorro mem. (\%)} & \textbf{Tiempo (s)} & \textbf{Tiempo rel. (x)} \\
        \hline
        22 & Referencia & 71.6 & 0.0 & 0.12 & 1.00 \\
        22 & Blosc & 15.8 & 77.9 & 140.33 & 1159.78 \\
        22 & ZFP & 38.9 & 45.7 & 1110.10 & 9174.34 \\
        23 & Referencia & 135.6 & 0.0 & 0.22 & 1.00 \\
        23 & Blosc & 16.1 & 88.1 & 277.28 & 1249.01 \\
        23 & ZFP & 64.0 & 52.8 & 3052.17 & 13748.49 \\
        24 & Referencia & 263.6 & 0.0 & 0.47 & 1.00 \\
        24 & Blosc & 17.0 & 93.6 & 558.54 & 1189.23 \\
        24 & ZFP & 113.2 & 57.1 & 4424.57 & 9420.66 \\
        \hline
    \end{tabular}
    }
\end{table}

\begin{table}[H]
    \centering
    \small
    \setlength{\tabcolsep}{3pt}
    \caption{Variante optimizada de Grover multiobjetivo}
    \label{tab:grover-optimized-multi}
    \resizebox{\linewidth}{!}{%
    \begin{tabular}{lccccc}
        \hline
        \textbf{Qubits} & \textbf{Estrategia} & \textbf{Mem. pico (MB)} & \textbf{Ahorro mem. (\%)} & \textbf{Tiempo (s)} & \textbf{Tiempo rel. (x)} \\
        \hline
        22 & Referencia & 71.7 & 0.0 & 0.10 & 1.00 \\
        22 & Blosc & 15.9 & 77.9 & 35.39 & 340.26 \\
        22 & ZFP & 39.0 & 45.6 & 279.05 & 2683.13 \\
        23 & Referencia & 135.6 & 0.0 & 0.26 & 1.00 \\
        23 & Blosc & 16.2 & 88.1 & 483.84 & 1890.01 \\
        23 & ZFP & 64.1 & 52.8 & 5298.98 & 20699.15 \\
        24 & Referencia & 263.7 & 0.0 & 0.47 & 1.00 \\
        24 & Blosc & 17.0 & 93.5 & 551.55 & 1173.51 \\
        24 & ZFP & 113.1 & 57.1 & 8305.50 & 17671.27 \\
        \hline
    \end{tabular}
    }
\end{table}

La tendencia cambia de forma importante respecto a la implementación principal. Dentro de la
variante optimizada, Blosc conserva ahorros muy altos de memoria, entre 77.9\% y 93.6\%, y
ZFP se sitúa entre 45.6\% y 57.1\%. Sin embargo, la referencia no comprimida de esta variante
ya ejecuta Grover en unas pocas décimas de segundo, por lo que cualquier compresión introduce
penalizaciones temporales extremas: entre $340.26\times$ y $1890.01\times$ para Blosc, y entre
$2683.13\times$ y $20699.15\times$ para ZFP en el caso multiobjetivo. En consecuencia, una vez
se explota la parametrización reducida de Grover, la compresión del estado deja de ser una
opción atractiva desde el punto de vista práctico, salvo en escenarios donde la memoria
disponible imponga una restricción mucho más severa que el tiempo de ejecución.
